\chapter{}

Toppa era morto di vecchiaia. Lo trovarono una mattina di febbraio, sotto il carro; nell’aia. Il gelo lo aveva attaccato mezzo ai mattoni; e la pancia, quando Carlo gli ci picchiò la pala che doveva adoprare per sotterrarlo a un olivo, suonò come un tamburo; e fece, perciò, ridere.

Era stato, dopo la castratura, piuttosto cattivo: quando non voleva esser toccato, prima si allontanava; e poi, se non smettevano, si avventava digrignando i denti. Era bastardo e alto un mezzo metro. Aveva quel pelo bianco che vicino alla pelle è giallo, con una macchia nera sopra un orecchio; e perciò gli trovarono quel nome.

Da piccolo, a pena slattato, Domenico lo legò al ferro del pozzo; e, quando guaiva, gli assalariati avevano l’ordine di pigliarlo a calci.

Poi gli comprò un collare con i chiodi d'ottone; un collare che non gli levavano mai altro che mentre lo tosavano.

Egli udiva la sonagliera del cavallo di Domenico quando ancora era al borgo fuor di Porta Camollia. Allora, esciva nella strada; e cominciava ad abbaiare. Quando il cavallo appariva ad una svolta poco distante dal cancello del podere, si metteva a correre da un punto all'altro della strada. Le persone si tiravano da parte; ma Toppa aveva buttato giù parecchi ragazzi, che non erano stati in tempo.

Quando aveva mangiato, andava invece a correre per i campi, e ci lasciava i segni da per tutto; specie dov'era il grano alto ci restava un solco che si vedeva anche di lontano. Quando seminavano, dovevano prenderlo a sassate perchè dove passava saltando bisognava rifare il lavoro. Gli piaceva l'uva matura e i fichi anche di più.

Obbediva soltanto a Domenico e a Giacco; degli altri aveva soltanto timore, quando non gli veniva voglia di mordere; come fece una volta a Ghìsola che gli era salita a cavallo.

Non c’era nessun altro cane che la potesse con lui; e ne fece morire più d’uno per averli azzannati su la spina dorsale. Due li sbranò perchè erano andati a mangiargli la zuppa nel catino.

Tollerava invece i gatti, purchè non gli andassero vicino. Ma quando stava al sole, non ce li voleva in nessun modo: teneva, allora, un occhio chiuso e un altro aperto: ne apriva uno e ne chiudeva un altro. All’improvviso, faceva un balzo con un abbaio che stordiva.

Non ebbe voglia di ruzzare nè meno da cucciolo. E si comportava a seconda di chi lo avvicinasse: non sbagliava. Non avrebbe obbedito a Pietro, nè mai gli fece una carezza.

Quando lo sotterrarono, dopo aver avvertito il Rosi, che ricordò di averlo pagato due lire soltanto, dando l’ordine di serbare il collare, Giacco pianse. Anch’egli si sentiva vecchio; e, guardando il cadavere della bestia, disse agli altri:

— Noi faremo la stessa fine.

Enrico rispose:

— Di più ormai non poteva campare. Che ci fanno i vecchi al mondo?

E dette un'occhiata a Carlo, che rideva.

Ma Giacco buttò via la zappa, e gridò:

— Io camperò più di te: mettitelo bene in mente. Vedi questa povera bestia? Aveva il cuore più buono del tuo!

— Io non ho voluto alludere a te.

— E a chi, dunque? Il cervello l'ho debole ora, ma la ragione l'ho sempre....

Carlo, allora, cominciò a bestemmiare e a pigliarsela con il cane:

— Non poteva campare? La fatica per la buca non ci sarebbe stata; e nè meno questa questione. Bada se per una carogna ci si deve offendere!

Egli fingeva d'essere arrabbiato; ma invece, aveva piacere che, senza compromettersi lui, Giacco facesse il viso bianco a quel modo. E Giacco guardava il cane, stando attento che gli altri non lo pestassero per sbadataggine e per dispetto.

Masa, venuta a vederlo metterlo sottoterra, si fermò un poco distante dalla buca; senza smettere di mangiare, sebbene si sentisse agitata. Quand'ebbero finito, si picchiò il ventre con un pugno, e disse:

— Se mangio dell'altro, le budella mi fanno gomìcciolo in corpo.

Giacco alzò la testa e la guardò:

— Vorrei ridere, allora! Piuttosto va' alle tue faccende. Creperesti prima di smettere! Lo capisci che mi fai rabbia?

Masa mise il pane in tasca, e rispose:

— Sei un gran brontolone! Il Signore lo sa!

Sospirò; e, seguitando a camminare innanzi agli altri, aggiunse tra sè:

— Pazienza, pazienza!

Ella non sapeva quel che avevano detto al marito.
